\chapter{Finite-time extinction}
\label{chap:finite-time-extinction}
\label{energy}

\section{The extinction theorem}

\begin{theorem}[Finite-time extinction for Ricci flow with surgery]
\label{thm:finite-time-extinction-main}
\uses{MAIN}
\label{extinct}

Let $(M,g(0))$ be a compact, connected, normalized Riemannian $3$-manifold
whose fundamental group is a free product of finite groups and infinite cyclic
groups. Then $M$ contains no embedded $\Ar P^2$ with trivial normal bundle. The
Ricci flow with surgery $({\mathcal M},G)$ supplied by
Theorem~\ref{MAIN}, with initial slice $(M,g(0))$ and defined for all
$t\in[0,\infty)$, is extinct in finite time: $M_T=\emptyset$ for all
sufficiently large $T$.
\end{theorem}
\begin{proof}
  \uses{prop:eventual-trivial-pi2, prop:w3-forward-difference-quotient, MAIN, lem:forward-difference-quotient, thm:finite-time-extinction-main, Tgood0good}

\sketch An embedded $\Ar P^2$ with trivial normal bundle can be made disjoint
from the connected-sum spheres. It lies in a prime factor, is nonseparating,
and induces an injection of its order-two fundamental group. Mayer--Vietoris
gives a surjection of the first homology of that factor onto $\Zee$, contrary
to the factor having either finite or infinite cyclic fundamental group; see
\cite{Hatcher,Hempel}. Theorem~\ref{MAIN} therefore supplies the flow.
\cite{AG,AG,AG,AG,AG,ColdingMinicozzi,HamiltonNSRF3M,Hempel,Hildebrandt,Jost,Perelman3,Perelman3,Perelman3,Perelman3,Perelman3,SacksUhlenbeck,SacksUhlenbeck,SacksUhlenbeck,SacksUhlenbeck,SchoenYau}.
Proposition~\ref{prop:eventual-trivial-pi2} eliminates nontrivial $\pi_2$ after
a finite time. Each surviving path of components then carries a nonzero
$\pi_3$ class by Claim~\ref{claim:x-t1-nontrivial-pi3}. Applying
Proposition~\ref{prop:w3-forward-difference-quotient} and
$R_{\min}(t)\ge-6/(1+4t)$ bounds its nonnegative width by the solution of
\[
w'=-2\pi+\frac{3}{1+4t}w.
\]
That solution becomes negative in finite time. There are finitely many
components at the initial comparison time, so all component paths terminate
under one uniform finite bound.
\end{proof}

\begin{corollary}[Topological conclusion of finite-time extinction]
Let $M$ satisfy the hypotheses of Theorem~\ref{Theorem1}. Then $M$ is a
connected sum of spherical space forms and $2$-sphere bundles over $S^1$. If
$M$ is simply connected, then $M\cong S^3$; if $\pi_1(M)$ is finite, then $M$
is a spherical space form.
\end{corollary}
\begin{proof}
\uses{thm:finite-time-extinction-main,MAIN,Tgood0good}
\sketch Use Theorem~\ref{MAIN} and
Theorem~\ref{thm:finite-time-extinction-main} to choose $T$ with
$M_T=\emptyset$. Corollary~\ref{Tgood0good} gives the stated connected-sum
description of $M=M_0$. The two special cases follow from the fundamental
group.
\end{proof}

\begin{definition}[Path of components]
\label{def:path-of-components}
Let $I$ be any interval and let $({\mathcal M},G)$ be a Ricci flow with
surgery over $I$. A \emph{path of components} is a connected open subset
${\mathcal X}\subset {\bf t}^{-1}(I)$ such that
${\mathcal X}(t)={\mathcal X}\cap M_t$ is a connected component of $M_t$ for
every $t\in I$.
\end{definition}

On an interval without an interior surgery time, a path of components is the
ordinary Ricci-flow continuation of any one of its slices. At a surgery time it
continues through one selected component of the surgered predecessor; it cannot
continue if that predecessor is discarded completely.

\section{Disappearance of components with nontrivial $\pi_2$}

\subsection{Essential sphere systems and surgeries}

\begin{lemma}[Rank bound for free product decompositions]
\label{lem:free-product-rank-bound}
\label{grouptheory}
If a group $G$ generated by $k$ elements has a decomposition
$G=G_1*\cdots*G_\ell$ with every $G_i$ nontrivial, then $\ell\le k$.
\end{lemma}
\begin{proof}
\sketch Grushko's theorem decomposes a rank-$k$ free group surjecting onto $G$
as $F_1*\cdots*F_\ell$, with $F_i\twoheadrightarrow G_i$; hence every $F_i$
has positive rank and $\ell\le k$ \cite{Stallings}.
\end{proof}

\begin{definition}[Homotopically essential family of 2-spheres]
\label{def:homotopically-essential-family}
For a compact, possibly disconnected $3$-manifold $X$, an embedded $2$-sphere
is \emph{homotopically essential} if its inclusion in $X$ is not null-homotopic.
A finite disjoint family $\{\Sigma_1,\ldots,\Sigma_n\}$ is homotopically
essential if each $\Sigma_i$ is essential and no two inclusion maps
$\Sigma_i\hookrightarrow X$ are homotopic.
\end{definition}

\begin{lemma}[Finiteness of homotopically essential families of 2-spheres]
\label{lem:homotopically-essential-family-finite}
For every compact $3$-manifold $X$, the cardinalities of its homotopically
essential disjoint sphere families have a finite upper bound.
\end{lemma}
\begin{proof}
\uses{lem:free-product-rank-bound}
\sketch Reduce to connected $X$. Cutting along such a family gives a graph-of-groups
decomposition of $\pi_1(X)$ with trivial edge groups \cite{Hatcher}. If $E$ is
the number of spheres, $V_i$ the number of vertices of valence $i$, and $k$ the
first Betti number of the underlying graph, then
\[
2V_1+V_2\ge E+3(1-k).
\]
The vertex groups together with $k$ infinite cyclic factors form a free-product
decomposition. Since $k\le\operatorname{rank}H_1(X)$, the preceding lemma and
finite generation of $\pi_1(X)$ bound $E$.
\end{proof}

\begin{definition}[Homotopically essential surgery]
\label{def:homotopically-essential-surgery}
A surgery along a sphere $S_0(t)$ at time $t$ is \emph{homotopically
essential} if, for every sufficiently close $t'<t$, the backward-flowed sphere
$S_0(t')\subset M_{t'}$ is homotopically essential.
\end{definition}

For compact $X$, let $s(X)$ be the largest size of a homotopically essential
disjoint sphere family, and set $s(t)=s(M_t)$.

\begin{claim}[Monotonicity of the essential-sphere invariant $s$]
\label{claim:invariant-s-monotone}
For $t'<t$, one has $s(t')\ge s(t)$. If surgery at $t$ uses at least one
homotopically essential sphere, then $s(t')>s(t)$ for all sufficiently close
$t'<t$.
\end{claim}
\begin{proof}
\uses{surgerytoptype}
\sketch By Proposition~\ref{surgerytoptype}, surgery consists of cutting along finitely
many spheres, capping, and discarding components. Trivial cuts preserve $s$ and
discarding components cannot increase it. Move an essential family in $M_t$
off the caps and pull it back to $M_{t'}$; adjoining an essential surgery sphere
gives a family with one more member. Null-homotopy or mutual homotopy in the
pulled-back family would either persist after surgery or make the surgery
sphere null-homotopic, both impossible.
\end{proof}

\begin{proposition}[Finitely many homotopically essential surgeries]
\label{prop:finitely-many-homotopically-essential-surgeries}
\label{homess}
Any Ricci flow with surgery satisfying Assumptions (1)--(7) of
Chapter~\ref{sect:surgery} has only finitely many homotopically essential
sphere surgeries.
\end{proposition}
\begin{proof}
  \uses{surgerytoptype}
\sketch By Lemma~\ref{lem:homotopically-essential-family-finite}, $s(M_0)$ is
finite. Claim~\ref{claim:invariant-s-monotone} makes $s(t)$ a nonnegative
integer that drops at every essential surgery.
\end{proof}

\subsection{The $2$-sphere width}

\begin{lemma}[Minimal energy gap for homotopically nontrivial 2-spheres]
\label{lem:minimal-sphere-energy-gap}
Let $X$ be a compact Riemannian manifold with $\pi_2(X)\ne0$. There is
$e_0=e_0(X)>0$ such that every map $S^2\to X$ of area below $e_0$ is
null-homotopic, while some non-null-homotopic branched minimal immersion
$f\colon S^2\to X$ has area $e_0$.
\end{lemma}
\begin{proof}
\sketch The positive energy gap is Theorem 3.3 of \cite{SacksUhlenbeck}.
Minimize
\[
E_\alpha(s)=\int_{S^2}(1+|ds|^2)^\alpha\,da,
\qquad \alpha>1,
\]
among nontrivial maps and let $\alpha\downarrow1$. Sacks--Uhlenbeck
compactness produces either a strong harmonic limit or a bubble. A trivial
bubble could be cut out and replaced by the complementary disk, lowering the
energy of a nontrivial representative by a fixed amount; this contradicts
minimality. A resulting nontrivial harmonic sphere realizes $e_0$, and
conformality makes it a branched immersion.
\end{proof}

For a path ${\mathcal X}$ whose slices have nontrivial $\pi_2$, write
$W_2(t)=e_0({\mathcal X}(t))$.

\begin{claim}[Variation of area of a minimal 2-sphere]
\label{claim:minimal-sphere-area-variation}
\label{varareaform}
If $f\colon S^2\to(M,g(t_0))$ is a branched minimal sphere, then
\[
\left.\frac d{dt}\Area_{g(t)}f(S^2)\right|_{t=t_0}
\le -4\pi-\frac12R_{\min}(g(t_0))\Area_{g(t_0)}f(S^2).
\]
\end{claim}
\begin{proof}
\sketch Under Ricci flow,
\begin{equation}\label{ha1}
\left.\frac d{dt}\Area_{g(t)}f(S^2)\right|_{t=t_0}
=-\int_{S^2}(R-\Ric({\bf n},{\bf n}))\,da.
\end{equation}
The Gauss equation and minimality give
\begin{equation}\label{ha2}
\left.\frac d{dt}\Area_{g(t)}f(S^2)\right|_{t=t_0}
=-\int_{S^2}K\,da-\frac12\int_{S^2}(|A|^2+R)\,da.
\end{equation}
The same identity holds off the finite branch set. Gauss--Bonnet yields
\begin{equation}\label{ha3}
\left.\frac d{dt}\Area_{g(t)}f(S^2)\right|_{t=t_0}
\le-4\pi-\frac12R_{\min}(t_0)\Area_{g(t_0)}f(S^2).
\end{equation}
\end{proof}

\begin{claim}[Area growth bound for a fixed 2-sphere map]
\label{claim:sphere-area-growth-bound}
Let $(M,g(t))$, $t_0\le t\le t_1$, be a compact Ricci flow with
$|\Ric_{g(t)}|\le D$. Every $C^1$ map $f\colon S^2\to M$ satisfies
\[
\Area_{g(t_1)}f(S^2)\le e^{4D(t_1-t_0)}\Area_{g(t_0)}f(S^2).
\]
\end{claim}
\begin{proof}
\sketch The metric variation gives
$\frac d{dt}\Area_{g(t)}f(S^2)\le4D\Area_{g(t)}f(S^2)$; integrate this
inequality.
\end{proof}

\begin{lemma}[Forward difference quotient for $W_2$]
\label{lem:w2-forward-difference-quotient}
\uses{lem:minimal-sphere-energy-gap, lem:forward-difference-quotient}
\label{W_2}

Along such a path of components,
\[
\frac{dW_2}{dt}(t)\le-4\pi-\frac12R_{\min}(t)W_2(t)
\]
in the sense of forward difference quotients. The function $W_2$ is continuous
away from surgery times, and at a surgery time $t$ it satisfies
\[
W_2(t)\le\liminf_{t'\to t^-}W_2(t').
\]
\end{lemma}
\begin{proof}
\uses{claim:minimal-sphere-area-variation}
\sketch At a regular time, use an area-minimizing essential sphere and
Claim~\ref{claim:minimal-sphere-area-variation}. Upper semicontinuity follows
by keeping a fixed sphere, while Claim~\ref{claim:sphere-area-growth-bound}
applied to minimizing spheres at nearby times gives lower semicontinuity. At a
surgery time, Proposition~\ref{surgerydist} supplies, for every $\eta>0$, a
$(1+\eta)$-Lipschitz homotopy equivalence from the preceding component to the
surviving one. Composing a minimizing sphere with this map and letting
$\eta\downarrow0$ proves the liminf inequality.
\end{proof}

\begin{proposition}[Eventual triviality of $\pi_2$]
\label{prop:eventual-trivial-pi2}
\label{pi2}
Let $({\mathcal M},G)$ satisfy the conclusion of Theorem~\ref{MAIN}. There is
$T_1<\infty$ such that every component of $M_T$ has trivial $\pi_2$ for all
$T\ge T_1$. Each such component has either finite fundamental group and a
homotopy $3$-sphere as universal cover, or infinite fundamental group and a
contractible universal cover.
\end{proposition}
\begin{proof}
\sketch Proposition~\ref{prop:finitely-many-homotopically-essential-surgeries}
gives a time $T_0$ after which every sphere surgery is homotopically trivial.
Each component of $M_{T_0}$ with nontrivial $\pi_2$ then has a unique maximal
continuation with nontrivial $\pi_2$. If one persisted forever,
Lemma~\ref{lem:w2-forward-difference-quotient}, the bound
$R_{\min}(t)\ge-6/(1+4t)$, and the comparison equation
\[
w_2'=-4\pi+\frac{3}{1+4t}w_2,
\qquad w_2(T_0)=W_2(T_0),
\]
would give
\[
w_2(t)=w_2(T_0)\frac{(4t+1)^{3/4}}{(4T_0+1)^{3/4}}
+4\pi(4T_0+1)^{1/4}(4t+1)^{3/4}-4\pi(4t+1),
\]
which is eventually negative, contradicting $W_2>0$. Finitely many initial
components give a common $T_1$. If a connected $3$-manifold has $\pi_2=0$,
its universal cover is $2$-connected; it is a homotopy sphere when compact,
and is contractible by Hurewicz and Whitehead when noncompact.
\end{proof}

\begin{corollary}[Sphere theorem for closed 3-manifolds without embedded $\mathbb{RP}^2$]
\label{cor:sphere-theorem-no-rp2}
\label{spherethm}
Let $M$ be a closed connected $3$-manifold with no embedded $\Ar P^2$ having
trivial normal bundle and with $\pi_2(M)\ne0$. Then either
$M=M_1\#M_2$, with neither $M_i$ homotopy equivalent to $S^3$, or $M$ has a
prime factor that is an $S^2$-bundle over $S^1$. In either case $M$ contains
an embedded homotopically nontrivial $2$-sphere.
\end{corollary}
\begin{proof}
\uses{prop:eventual-trivial-pi2}
\sketch Run the all-time surgery flow from a normalized metric and take the first
surgery that kills $\pi_2$. It either cuts along an essential sphere or removes
an $S^2$-bundle over $S^1$ or $\Ar P^3\#\Ar P^3$. In each case an essential
embedded sphere pulls back through all earlier, homotopically trivial surgeries
to $M$.
\end{proof}

\begin{remark}[Components with nontrivial $\pi_2$ that disappear]
\label{rem:pi2-disappearing-components}
\uses{thm:finite-time-extinction-main}

The discarded components with nontrivial $\pi_2$ are exactly the
$S^2\times\Ar$-type components: $S^2$-bundles over $S^1$ and
$\Ar P^3\#\Ar P^3$. Consequently, after essential sphere surgeries cease,
every other component has universal cover either contractible or a homotopy
$3$-sphere.
\end{remark}

\section{The $\pi_3$ width}

Let $\Lambda M$ be the free loop space of $C^1$ maps $S^1\to M$ with the
$C^1$ topology, and let $*$ denote the constant loop at a fixed $x_0\in M$.

\begin{claim}[$\pi_2$ of the loop space identified with $\pi_3$]
\label{claim:loop-space-pi2-pi3-iso}
If $\pi_2(M,x_0)=0$, then
$\pi_2(\Lambda M,*)\cong\pi_3(M,x_0)$. It is represented by free homotopy
classes of maps from $S^2$ into the component of null-homotopic loops.
\end{claim}
\begin{proof}
\sketch A based map $S^2\to\Lambda M$ is a map $S^2\times S^1\to M$ that
collapses $\{\mathrm{pt}\}\times S^1$. The quotient is homotopy equivalent to
$S^2\vee S^3$, so its based homotopy classes are
$\pi_2(M)\oplus\pi_3(M)$. The first summand vanishes. The free description
follows by lifting to the universal cover and identifying $\pi_3(M)$ with its
$H_3$ by Hurewicz.
\end{proof}

\begin{definition}[Spanning-disk energy functional $W(\xi)$]
\label{def:w-xi-spanning-disk-functional}
\uses{prop:eventual-trivial-pi2, lem:w2-forward-difference-quotient}
\label{Wdefn}

For a null-homotopic loop $\gamma\in\Lambda M$, let $A(\gamma)$ be the
infimum of areas of Lipschitz spanning disks, with boundary equal to $\gamma$
up to reparameterization. For a continuous family
$\Gamma\colon S^2\to\Lambda M$ of null-homotopic loops, set
\[
W(\Gamma)=\max_{c\in S^2}A(\Gamma(c)).
\]
For $\xi\in\pi_2(\Lambda M,*)$, let $W(\xi)$ be the infimum of $W(\Gamma)$
over all, not necessarily based, representatives of $\xi$ in the null-loop
component.
\end{definition}

A nonzero $\xi(t_0)\in\pi_3({\mathcal X}(t_0))$ is transported along a path
of components by Ricci-flow identifications and, at surgery, by the
distance-decreasing degree-one map of Proposition~\ref{surgerydist}. When
$\pi_2({\mathcal X}(t))=0$, write $W_\xi(t)=W(\xi(t))$.

\begin{claim}[Continuity of $W_\xi$ away from surgery times]
\label{claim:w-xi-continuity-nonsurgery}
At every nonsurgery time $t$, the function $W_\xi$ is continuous.
\end{claim}
\begin{proof}
\sketch Nearby slices are canonically identified. A fixed family of spanning
disks gives upper semicontinuity, while comparison of the nearby metrics by a
factor $1+\eta(|t-t'|)$ gives lower semicontinuity.
\end{proof}

\begin{claim}[Upper semicontinuity of $W_\xi$ at surgery times]
\label{claim:w-xi-surgery-semicontinuity}
At a surgery time $t$,
\[
W_\xi(t)\le\liminf_{t'\to t^-}W_\xi(t').
\]
\end{claim}
\begin{proof}
\sketch Proposition~\ref{surgerydist} gives, for close $t'<t$, a
$(1+\eta)$-Lipschitz map ${\mathcal X}(t')\to{\mathcal X}(t)$ carrying the
class $\xi(t')$ to $\xi(t)$. Apply it to a nearly minimizing disk family and
let $\eta\downarrow0$.
\end{proof}

\begin{definition}[Comparison ODE solution $w_{a,t'}$]
\label{def:w-a-t-ode-solution}
For a compact Ricci flow $(M,g(t))$, $t_0\le t\le t_1$, let $w_{a,t'}$ solve
\begin{equation}\label{wdiffeqn}
\frac{dw_{a,t'}}{dt}=-2\pi-\frac12R_{\min}(t)w_{a,t'}(t),
\qquad w_{a,t'}(t')=a.
\end{equation}
Write $w_a=w_{a,t_0}$.
\end{definition}

\begin{definition}[Auxiliary integral $A(t)$]
\label{def:a-function-integral}
For fixed $t'$, set
\[
A(t)=\int_{t'}^t\frac12R_{\min}(s)\,ds.
\]
\end{definition}

\begin{claim}[Explicit formula for $w_{a,t'}$]
\label{claim:w-a-t-explicit-formula}
\label{vclaim}
For $t''\ge t'$,
\[
w_{a,t'}(t'')=e^{-A(t'')}
\left(a-2\pi\int_{t'}^{t''}e^{A(t)}\,dt\right).
\]
If $a'>a$, then
\[
w_{a',t'}(t'')=w_{a,t'}(t'')+(a'-a)e^{-A(t'')}.
\]
\end{claim}
\begin{proof}
\sketch Multiply Equation~\eqref{wdiffeqn} by $e^{A(t)}$ and integrate; the
second identity is the difference of two solutions.
\end{proof}

\begin{lemma}[Families of short loops represent the trivial $\pi_3$ class]
\label{lem:short-loops-trivial-pi3}
\label{shortpi2triv}
Let $(X,g)$ be compact with $\pi_2(X)=0$. There is $\zeta>0$ such that every
family $\Gamma\colon S^2\to\Lambda X$ whose loops all have length below
$\zeta$ represents zero in $\pi_3(X)$.
\end{lemma}
\begin{proof}
\sketch Choose $\zeta$ below the injectivity radius. Join each point of each
loop to its value at a fixed parameter by the unique short geodesic. These
geodesics give a continuous family of spanning disks and deform $\Gamma$ to a
family of constant loops. That remaining family is null-homotopic because
$\pi_2(X)=0$.
\end{proof}

\begin{corollary}[Short loops bound small-area disks]
\label{cor:short-loop-bounds-small-disk}
\uses{lem:short-loops-trivial-pi3, claim:w-a-t-explicit-formula}
\label{shlnsharea}

For every compact Riemannian manifold $(X,g)$ and every $\eta>0$, there is
$0<\zeta<\eta/2$ such that every $C^1$ loop of length less than $\eta$ bounds
a disk of area less than $\eta$.
\end{corollary}
\begin{proof}
\sketch Use radial geodesics inside a normal ball; the cone area tends
uniformly to zero with the loop length.
\end{proof}

\begin{proposition}[Forward difference quotient for $W_\xi$]
\label{prop:w3-forward-difference-quotient}
\uses{prop:eventual-trivial-pi2, def:w-xi-spanning-disk-functional}
\label{W_3fordiff}

Let $({\mathcal M},G)$ be as in Theorem~\ref{MAIN}, let ${\mathcal X}$ be a
path of components on $[t_0,t_1]$ with
$\pi_2({\mathcal X}(t_0))=0$, and let
$0\ne\xi\in\pi_3({\mathcal X}(t_0),*)$. Then
\[
\frac{dW_\xi}{dt}(t)\le-2\pi-\frac12R_{\min}(t)W_\xi(t)
\]
in the sense of forward difference quotients. The function is continuous at
nonsurgery times and satisfies
$W_\xi(t)\le\liminf_{t'\to t^-}W_\xi(t')$ at surgery times.
\end{proposition}
\begin{proof}
  \uses{prop:eventual-trivial-pi2, prop:w3-forward-difference-quotient, MAIN, lem:forward-difference-quotient}
\sketch Continuity and surgery semicontinuity are
Claims~\ref{claim:w-xi-continuity-nonsurgery} and
\ref{claim:w-xi-surgery-semicontinuity}. On a surgery-free interval, choose a
family within $\varepsilon$ of $W_\xi(t_0)$ and apply
Proposition~\ref{prop:curve-deformation-existence}. If the comparison solution
is nonnegative, Corollary~\ref{cor:short-loop-bounds-small-disk} controls the
short-loop alternative and gives $W_\xi(t_1)\le w(t_1)+o(1)$. If it is
negative, all non-short alternatives would have negative area, so every loop
would be short, contradicting Lemma~\ref{lem:short-loops-trivial-pi3}. Letting
$t_1\downarrow t_0$ gives the quotient inequality.
\end{proof}

\begin{claim}[Nontrivial free product or $\mathbb{Z}$ fundamental group forces $\pi_2\ne0$]
\label{claim:free-product-implies-pi2-nonzero}
If $X$ is a compact $3$-manifold and $\pi_1(X)$ is either a nontrivial free
product or isomorphic to $\Zee$, then $\pi_2(X)\ne0$.
\end{claim}
\begin{proof}
\sketch This is the splitting result for compact $3$-manifolds; see
Theorems 5.2 and 7.1 of \cite{Hempel}.
\end{proof}

\begin{claim}[$\mathcal{X}(T_1)$ has nontrivial $\pi_3$]
\label{claim:x-t1-nontrivial-pi3}
Under the hypotheses of Theorem~\ref{thm:finite-time-extinction-main}, every
component ${\mathcal X}(T_1)$ surviving after the time in
Proposition~\ref{prop:eventual-trivial-pi2} has nontrivial $\pi_3$.
\end{claim}
\begin{proof}
\sketch The surgery topology preserves the property that each component group
is a free product of finite and infinite cyclic groups. Since
$\pi_2({\mathcal X}(T_1))=0$, Claim~\ref{claim:free-product-implies-pi2-nonzero}
forces $\pi_1({\mathcal X}(T_1))$ to be finite. Its universal cover is a closed
simply connected $3$-manifold, hence a homotopy $3$-sphere by Poincare duality,
so $\pi_3({\mathcal X}(T_1))\cong\Zee$.
\end{proof}

\section{Curve-shrinking flow}
\label{sect:csf}

For a Ricci flow $(M,g(t))$ and an immersed family
$c\colon S^1\times[t_0,t_1]\to M$, set
\[
X=\partial_xc,\qquad S=\frac X{|X|},\qquad
H=\nabla^g_SS,\qquad k=|H|.
\]
The curve-shrinking equation is $\partial_tc=H$. Arc-length derivatives are
denoted by a prime.

\begin{claim}[Short-time existence of the curve-shrinking flow]
\label{claim:curve-shrinking-short-time-existence}
\uses{prop:curve-deformation-existence}
\label{curvshrexist}

Every immersed $C^2$ loop at time $t_0$ has an immersed curve-shrinking flow
on $[t_0,t_1')$ for some $t_1'>t_0$. The solution extends through $t_1'$ as
immersions unless $\max_{S^1}k(\cdot,t)\to\infty$ as $t\uparrow t_1'$.
\end{claim}
\begin{proof}
\sketch This is the standard short-time existence and continuation criterion,
Theorem 1.13 of \cite{AG}.
\end{proof}

\begin{lemma}[Area bound for embedded curve-shrinking flow]
\label{lem:embedded-curve-shrinking-area-bound}
\label{embcase}
Let $c(\cdot,t)$, $t_0\le t\le t_1$, be an embedded curve-shrinking flow of
null-homotopic loops, and let $A(t)$ be the least spanning-disk area. Then
$A$ is continuous and
\[
\frac{dA}{dt}(t)\le-2\pi-\frac12R_{\min}(t)A(t)
\]
in the sense of forward difference quotients.
\end{lemma}
\begin{proof}
\uses{claim:minimal-sphere-area-variation}
\sketch An area-minimizing spanning disk exists and is immersed
\cite{Hildebrandt,Morrey,HardtSimon,GL}. Extend the boundary motion to an
ambient isotopy. Metric variation, minimality, and the boundary velocity give
\[
\frac d{dt}\Area(D)\le-\int_DK\,da-\int_{\partial D}k_{\mathrm{geod}}\,ds
-\frac12R_{\min}\Area(D).
\]
Gauss--Bonnet turns the first two terms into $-2\pi$.
\end{proof}

\begin{corollary}[Area comparison for embedded curve-shrinking flow]
\label{cor:embedded-curve-shrinking-area-comparison}
Under the preceding hypotheses on $[t^-,t^+]$, let
$\psi'=-2\pi-\frac12R_{\min}\psi$ with $\psi(t^-)=A(t^-)$. Then
$A(t^+)\le\psi(t^+)$.
\end{corollary}
\begin{proof}
\sketch Apply Lemma~\ref{lem:forward-difference-quotient} to
Lemma~\ref{lem:embedded-curve-shrinking-area-bound}.
\end{proof}

\begin{claim}[Reduction to a $C^2$-approximating family]
\label{claim:curve-shrinking-approximation-suffices}
To prove the same area comparison for an immersed flow in dimension at least
$3$, it suffices, for every $\delta>0$, to find a $C^2$ family $\hat c$ within
$\delta$ of the original flow such that
\[
A(\hat c(t^+))\le\psi_{\delta,\hat c}(t^+),
\quad
\psi_{\delta,\hat c}'=-2\pi+2\delta L_{\hat c}
-\frac12R_{\min}\psi_{\delta,\hat c},
\]
with $\psi_{\delta,\hat c}(t^-)=A(\hat c(t^-))$.
\end{claim}
\begin{proof}
\sketch Let $\delta\downarrow0$. Continuity of least spanning area in the
$C^1$ topology and uniform convergence of the comparison ODEs yield the
unperturbed inequality.
\end{proof}

\begin{lemma}[Area bound for immersed curve-shrinking flow in dimension at least 3]
\label{lem:immersed-curve-shrinking-area-bound}
\label{nonembeasycase}
Let $\dim M\ge3$ and let $c(\cdot,t)$ be a $C^2$ curve-shrinking flow of
null-homotopic immersed loops on $[t^-,t^+]$. Its least spanning area $A(t)$
is continuous and satisfies $A(t^+)\le\psi(t^+)$ for the comparison solution
of the preceding corollary.
\end{lemma}
\begin{proof}
\uses{lem:embedded-curve-shrinking-area-bound,lem:forward-difference-quotient}
\sketch Approximate the family generically in $C^2$. In dimension at least $3$ the
approximating curves are embedded except at finitely many times. On every
embedded subinterval the boundary-variation error is bounded by
$2\delta L_{\hat c}$, giving the ODE in
Claim~\ref{claim:curve-shrinking-approximation-suffices}. Add the interval
estimates and send $\delta\to0$.
\end{proof}

\subsection{Space-time evolution estimates}

\begin{notation}[Space-time metric for the curve-shrinking flow]
\label{not:curve-shrinking-spacetime}
\label{spacetimehat}
On $M\times I$, $I=[t_0,t_1]$, let $\hat g=g(t)\oplus dt^2$ and let $\nabla$
and $\widehat{\Rm}$ denote its connection and curvature. Write $\nabla^g$ and
$\Ric_g$ for the horizontal connection and Ricci tensor. For a
curve-shrinking flow define
\[
\hat c(x,t)=(c(x,t),t),\qquad \widehat H=\hat c_*(\partial_t)=H+\partial_t.
\]
\end{notation}

\begin{lemma}[Christoffel symbols of the space-time metric]
\label{lem:spacetime-christoffel}
\uses{not:curve-shrinking-spacetime}
\label{hatchristoffel}

For horizontal indices $i,j,k$ and time index $0$,
\[
\widehat\Gamma^k_{ij}=\Gamma^k_{ij},\qquad
\widehat\Gamma^k_{i0}=\widehat\Gamma^k_{0i}=-g^{k\ell}(\Ric_g)_{i\ell},
\qquad \widehat\Gamma^0_{ij}=(\Ric_g)_{ij},
\]
and all remaining symbols vanish. Hence, for horizontal $A,B$,
\begin{equation}\label{hatEq1}
\nabla_AB=\nabla^g_AB+\Ric_g(A,B)\partial_t,
\end{equation}
\begin{equation}\label{hatEq2}
\langle\nabla_A\partial_t,B\rangle=-\Ric_g(A,B).
\end{equation}
\end{lemma}
\begin{proof}
\uses{not:curve-shrinking-spacetime}
\sketch Insert $\partial_tg=-2\Ric_g$, $\hat g_{00}=1$, and $\hat g_{i0}=0$ into the
Koszul formula.
\end{proof}

\begin{lemma}[First variation formulas for curve-shrinking flow]
\label{lem:curve-shrinking-first-variation}
\source{morgan-tian-correction}
\uses{not:curve-shrinking-spacetime, lem:spacetime-christoffel}
\label{1stcurvshr}


With the preceding notation,
\[
\widehat H(|X|^2)=-2\Ric_g(X,X)-2k^2|X|^2,
\qquad
[\widehat H,S]=(k^2+\Ric_g(S,S))S.
\]
\end{lemma}
\begin{proof}
\uses{lem:spacetime-christoffel}
\sketch Since $[\widehat H,X]=0$, differentiate $|X|^2$ using
Equation~\eqref{hatEq2}. Orthogonality of $H$ and $S$ gives the $-2k^2|X|^2$
term. Differentiate $S=|X|^{-1}X$ and substitute the first identity.
\end{proof}

\begin{lemma}[Evolution equation for the curvature $k^2$]
\label{lem:curve-shrinking-curvature-evolution}
\source{morgan-tian-correction}
\uses{not:curve-shrinking-spacetime, lem:spacetime-christoffel}
\label{2ndcurvshr}


With primes denoting arc-length derivatives,
\begin{align*}
\partial_tk^2={}&(k^2)''-2\langle(\nabla_SH)^\perp,(\nabla_SH)^\perp\rangle+2k^4
-2\Ric_g(H,H)+4k^2\Ric_g(S,S)\\
&+2\widehat{\Rm}(\widehat H,S,H,S)
+2\Ric_g(S,S)\Ric_g(H,H)-2(\nabla^g_S\Ric_g)(S,H).
\end{align*}
\end{lemma}
\begin{proof}
\uses{lem:curve-shrinking-first-variation,lem:spacetime-christoffel}
\sketch Differentiate $k^2=|H|^2$ in the $\widehat H$ direction, commute $\widehat H$
past $S$, and use Equations~\eqref{hatEq1}--\eqref{hatEq2}. The second
$S$-derivative splits as
\begin{equation}\label{hatEq3}
2\langle\nabla_S\nabla_S\widehat H,H\rangle
=2\langle\nabla_S\nabla_S\partial_t,H\rangle
+2\langle\nabla_S\nabla_SH,H\rangle.
\end{equation}
The first term supplies $-2(\nabla^g_S\Ric_g)(S,H)-2\Ric_g(H,H)$; the second
supplies $(k^2)''-2|(\nabla_SH)^\perp|^2-2k^4$. Substituting
$[\widehat H,S]=(k^2+\Ric_g(S,S))S$ gives the displayed formula.
\end{proof}

\begin{remark}[The book's Lemma 19.7 is not correct]
\label{rem:curve-shrinking-evolution-erratum}
\source{morgan-tian-correction}

A computation in a frozen pullback metric omits the last two terms of
Lemma~\ref{lem:curve-shrinking-curvature-evolution} and replaces
$2\widehat{\Rm}(\widehat H,S,H,S)$ by $\Rm(H,S,H,S)$. The omitted connection
variation forces the linear $k$ term below and, consequently, dependence on
initial length in the total-curvature estimate.
\end{remark}

\begin{corollary}[Differential inequality for $k^2$]
\label{cor:curve-shrinking-curvature-inequality-bound}
\source{morgan-tian-correction}
\uses{lem:curve-shrinking-curvature-evolution}
\label{*4'}


There is $\widehat C<\infty$, depending only on bounds for the ambient
curvature and its first covariant derivative, such that
\[
\partial_tk^2\le(k^2)''-2|(\nabla_SH)^\perp|^2+2k^4
+\widehat C(k^2+k).
\]
\end{corollary}
\begin{proof}
\uses{lem:curve-shrinking-curvature-evolution}
\sketch Bound the five ambient terms in
Lemma~\ref{lem:curve-shrinking-curvature-evolution}; terms with two factors of
$H$ are $O(k^2)$, while the mixed space-time curvature and $\nabla\Ric$ terms
also contribute $O(k)$.
\end{proof}

\begin{claim}[Differential inequality for the curvature $k$]
\label{claim:curve-shrinking-curvature-inequality}
\source{morgan-tian-correction}
\uses{cor:curve-shrinking-curvature-inequality-bound}
\label{3rdcurvshr}


For $h_\epsilon=(k^2+\epsilon^2)^{1/2}$ there is $C_1$, depending only on the
ambient flow, such that
\[
\partial_th_\epsilon\le h_\epsilon''+k^3+C_1(h_\epsilon+1).
\]
\end{claim}
\begin{proof}
\uses{cor:curve-shrinking-curvature-inequality-bound}
\sketch Differentiate $h_\epsilon^2$. The bound
$|(h_\epsilon)'|^2\le|(\nabla_SH)^\perp|^2$ absorbs the negative gradient term
in the preceding corollary. Divide by $2h_\epsilon$ and use
$k^4/h_\epsilon\le k^3$ and $(k^2+k)/h_\epsilon\le h_\epsilon+1$.
\end{proof}

Set
\[
L(t)=\int ds,
\qquad
\Theta(t)=\int k\,ds.
\]

\begin{lemma}[Growth bounds for length and total curvature]
\label{lem:curve-shrinking-length-curvature-growth}
\source{morgan-tian-correction}
\uses{claim:curve-shrinking-curvature-inequality}
\label{1835}


If $C_2$ bounds $|\Ric_g|$ and $C_1$ is as above, then
\begin{equation}\label{Lformula}
\frac{dL}{dt}\le\int(C_2-k^2)\,ds
\end{equation}
and
\[
\frac{d\Theta}{dt}\le(C_1+C_2)\Theta+C_1L.
\]
\end{lemma}
\begin{proof}
\uses{lem:curve-shrinking-first-variation,claim:curve-shrinking-curvature-inequality}
\sketch The first variation gives
\begin{equation}\label{dLdteqn}
\frac{dL}{dt}=\int(-\Ric_g(S,S)-k^2)\,ds.
\end{equation}
For $\Theta_\epsilon=\int h_\epsilon ds$, combine
Claim~\ref{claim:curve-shrinking-curvature-inequality} with the length-element
variation. The $k^3$ term is canceled by $-h_\epsilon k^2$, and the integral
of $h_\epsilon''$ vanishes. Let $\epsilon\downarrow0$.
\end{proof}

\begin{corollary}[Exponential bound for length and total curvature]
\label{cor:curve-shrinking-length-curvature-exponential-bound}
\source{morgan-tian-correction}
\uses{lem:curve-shrinking-length-curvature-growth}
\label{lncurvgrowth}


For $t'\le t''$,
\[
L(t'')\le L(t')e^{C_2(t''-t')},
\qquad
\Theta(t'')+L(t'')\le(\Theta(t')+L(t'))e^{(C_1+C_2)(t''-t')}.
\]
In particular,
$\Theta(t'')\le(\Theta(t')+L(t'))e^{(C_1+C_2)(t''-t')}$.
\end{corollary}
\begin{proof}
\uses{lem:curve-shrinking-length-curvature-growth}
\sketch Apply Gronwall first to $L$ and then to $L+\Theta$.
\end{proof}

\begin{remark}[The bound on total curvature depends on the initial length]
\label{rem:total-curvature-needs-initial-length}
\source{morgan-tian-correction}

The term $C_1L$ prevents control of $\Theta(t'')$ by $\Theta(t')$ alone. The
later compact families have uniform initial bounds for both quantities.
\end{remark}

\subsection{Ramps in $M\times S^1$}

Let $S^1_\lambda$ have length $\lambda$, let $U$ be its unit tangent field,
and for a curve in $M\times S^1_\lambda$ set $u=\langle S,U\rangle$.

\begin{claim}[Evolution equation for the ramp function $u$]
\label{claim:ramp-function-evolution}
\label{rampcomp}
Along curve-shrinking flow,
\[
\partial_tu=u''+(k^2+\Ric(S,S))u\ge u''-C'u,
\]
where $C'$ bounds the ambient Ricci norm.
\end{claim}
\begin{proof}
\uses{lem:curve-shrinking-first-variation}
\sketch The product field $U$ is parallel and $\Ric(V,U)=0$. Differentiate
$u=\langle S,U\rangle$, use the commutator formula for $S$, and identify
$\langle\nabla_SH,U\rangle=u''$.
\end{proof}

\begin{definition}[Ramp]
\label{def:ramp-curve}
A curve in $M\times S^1_\lambda$ is a \emph{ramp} if
$u=\langle S,U\rangle>0$ everywhere.
\end{definition}

\begin{corollary}[Ramps remain ramps under curve-shrinking flow]
\label{cor:ramp-preserved-under-flow}
\label{ubd}
If a curve-shrinking flow starts from a ramp, every later immersed curve is a
ramp.
\end{corollary}
\begin{proof}
\uses{claim:ramp-function-evolution}
\sketch The maximum principle applied to $e^{C't}u$ preserves a positive lower bound.
\end{proof}

\begin{lemma}[Curve-shrinking flow exists for all time for ramps]
\label{lem:curve-shrinking-exists-for-ramps}
\label{csforramps}
Let $(M,g(t))$, $t_0\le t\le t_1$, be a Ricci flow. Every ramp in
$M\times S^1_\lambda$ has a curve-shrinking flow on all of $[t_0,t_1]$, and
all its slices are ramps.
\end{lemma}
\begin{proof}
\uses{claim:curve-shrinking-short-time-existence,claim:curve-shrinking-curvature-inequality,claim:ramp-function-evolution,cor:ramp-preserved-under-flow}
\sketch \source{morgan-tian-correction}
For $h=(k^2+\epsilon^2)^{1/2}$, combine the inequalities for $h$ and $u$ to
obtain
\[
\partial_t(h/u)\le(h/u)''+2(u'/u)(h/u)'+C'(h+1)/u.
\]
The maximum principle bounds $h/u$ exponentially. Since $u$ stays uniformly
positive on finite intervals, $k$ remains bounded, so the short-time
continuation criterion extends the flow to $t_1$.
\end{proof}

\begin{lemma}[Growth rate of the minimal annulus between ramps]
\label{lem:minimal-annulus-area-growth}
\label{mugrowth}
Let $\dim M=n\ge3$, and let $c_1,c_2$ be degree-one ramp flows in
$M\times S^1_\lambda$. If $\mu(t)$ is the least area of an annulus between
them, then $\mu$ is continuous and
\[
\frac{d\mu}{dt}(t)\le(2n-1)\max_M|\Rm(\cdot,t)|\,\mu(t)
\]
in the forward-difference sense.
\end{lemma}
\begin{proof}
\sketch A least annulus exists and is regular away from finitely many branch
points \cite{Hildebrandt,Morrey}. After analytic approximation
\cite{Bando,GageHamilton}, its first variation is
\[
\frac d{dt}\Area(A)=\int_A-\operatorname{Tr}(\Ric^T)\,da
--\int_{\partial A}k_{\mathrm{geod}}\,ds.
\]
Gauss--Bonnet, including the nonpositive branch correction, and minimality
bound this by $(2n-1)\max|\Rm|\Area(A)$. Approximate intersecting boundary
curves by a third disjoint curve and pass to the limit.
\end{proof}

\begin{corollary}[Exponential growth bound for the minimal annulus area]
\label{cor:minimal-annulus-exponential-growth}
\uses{prop:curve-deformation-existence}

The least annulus area between two degree-one ramp flows grows at most
exponentially, with exponent controlled by the ambient sectional-curvature
bound and independent of $\lambda$.
\end{corollary}
\begin{proof}
\sketch Integrate Lemma~\ref{lem:minimal-annulus-area-growth}.
\end{proof}

\section{Approximation and deformation of loop families}

For a loop $c$ in $M$, define its degree-one lift to
$M\times S^1_\lambda$ by $c^\lambda(x)=(c(x),x)$.

\begin{definition}[Regular n-polygonal approximation]
\label{def:n-polygonal-approximation}
For $\xi_n=e^{2\pi i/n}$ and a $C^1$ loop $c$, the \emph{regular
$n$-polygonal approximation} $c_n$ joins
$c(\xi_n^k)$ to $c(\xi_n^{k+1})$ by a minimizing geodesic, parameterized at
constant speed, for $k=1,\ldots,n$.
\end{definition}

\begin{claim}[Basic properties of n-polygonal approximations]
\label{claim:polygonal-approximation-properties}
\label{onec}
For every $C^1$ loop $c$ and $\zeta>0$, all sufficiently fine regular
polygonal approximations $c_n$ have length within $\zeta$ of $L(c)$ and are
joined to $c$ by a piecewise smooth annulus of area below $\zeta$.
\end{claim}
\begin{proof}
\sketch Polygonal lengths converge to $L(c)$. For large $n$, $c_n$ is inside a
normal neighborhood of $c$; the pointwise short geodesics between them sweep
out annuli whose areas tend to zero.
\end{proof}

\begin{claim}[Uniformity of the polygonal approximation over a compact family]
\label{claim:polygonal-approximation-uniform}
\uses{claim:polygonal-approximation-properties}
\label{uniformc}

If $X\subset\Lambda M$ is compact, the threshold in the preceding claim may
be chosen uniformly for $c\in X$.
\end{claim}
\begin{proof}
\uses{claim:polygonal-approximation-properties}
\sketch Otherwise choose a violating sequence with $n\to\infty$; compactness gives a
limit loop, and the estimates for that loop persist in a neighborhood.
\end{proof}

\begin{corollary}[Polygonal approximation of a family $\Gamma$]
\label{cor:polygonal-approximation-of-family}
\label{Gammaapprox}
Let $\Gamma\colon S^2\to\Lambda M$ be a continuous family of null-homotopic
loops. For every $\zeta>0$ and all sufficiently large $n$, the regular
polygonal family $\Gamma_n$ is continuous; every $\Gamma_n(c)$ is
null-homotopic, has length within $\zeta$ of $L(\Gamma(c))$, and satisfies
\[
|A(\Gamma_n(c))-A(\Gamma(c))|<\zeta.
\]
\end{corollary}
\begin{proof}
  \uses{claim:polygonal-approximation-uniform, cor:polygonal-approximation-of-family}

\sketch Uniformly short sides are unique and vary continuously with their endpoints.
The small connecting annuli preserve homotopy classes and change least disk
area by at most their area.
\end{proof}

Fix a smooth periodic $\psi_n\ge0$, vanishing to infinite order at the
$n$th roots of unity, symmetric and monotone on each half-edge, and normalized
by $\int_1^{\xi_n}\psi_n=2\pi/n$. The induced homeomorphism
$\widetilde\psi_n(x)=\int_1^x\psi_n$ smooths a polygonal loop by
$\widetilde c_n=c_n\circ\widetilde\psi_n$ without changing its length or
least spanning area.

\begin{claim}[Length and curvature bounds for the ramp of the smoothed approximation]
\label{claim:smoothed-approximation-ramp-bounds}
For $0<\lambda<1$, the lift $\widetilde\Gamma(c)^\lambda$ has length at most
$\lambda+L(\Gamma(c))$ and total curvature at most $n\pi$.
\end{claim}
\begin{proof}
\sketch On each edge the tangent is
$L\psi_n(x)u+(\lambda/2\pi)v$ in a flat totally geodesic product strip. Its
length element is at most $(L\psi_n+\lambda/2\pi)dx$, and its direction turns
by less than $\pi/2$ on each half-edge. Sum over the $n$ edges.
\end{proof}

\begin{lemma}[Good C2-approximation of a family of loops]
\label{lem:good-approximation-family}
\label{goodapprox}
For a continuous $\Gamma\colon S^2\to\Lambda M$ representing a class in
$\pi_3(M,*)$ and $0<\zeta<1$, there is a continuous
$\widetilde\Gamma\colon S^2\to\Lambda M$ such that:
\begin{enumerate}
\item $[\widetilde\Gamma]=[\Gamma]$;
\item every $\widetilde\Gamma(c)$ is $C^2$;
\item its length differs from that of $\Gamma(c)$ by less than $\zeta$;
\item $|A(\widetilde\Gamma(c))-A(\Gamma(c))|<\zeta$;
\item for all $c$ and $0<\lambda<1$, both the length and total curvature of
$\widetilde\Gamma(c)^\lambda$ are bounded by a constant $C_0$ depending only
on $\Gamma$, $\zeta$, and the ambient curvature bound.
\end{enumerate}
\end{lemma}
\begin{proof}
\uses{cor:polygonal-approximation-of-family}
\sketch Choose one sufficiently fine polygonal approximation uniformly over the
family, then apply the fixed smoothing above. The first four properties follow
from the small annuli in
Corollary~\ref{cor:polygonal-approximation-of-family}; the last is
Claim~\ref{claim:smoothed-approximation-ramp-bounds}.
\end{proof}

For $c\in S^2$, let $\widetilde\Gamma_c^\lambda(t)$ be the all-time ramp flow
from $\widetilde\Gamma(c)^\lambda$, and let $p_1$ be projection to $M$.

\begin{claim}[Linear bound on the forward difference quotient of area]
\label{claim:area-forward-difference-linear-bound}
\label{areachange}
There is $C_4$, depending only on $t_1-t_0$, the ambient sectional-curvature
bound, $\Gamma$, and $\zeta$, such that for all $c$, $\lambda$, and
$t_0\le t'<t''\le t_1$,
\[
A(p_1\widetilde\Gamma_c^\lambda(t''))
-A(p_1\widetilde\Gamma_c^\lambda(t'))\le C_4(t''-t').
\]
\end{claim}
\begin{proof}
\uses{lem:curve-shrinking-length-curvature-growth}
\sketch The surface swept out between $t'$ and $t''$ has area growth controlled by its
ambient Ricci variation and by $\int k\,ds$, uniformly bounded by the length
and total-curvature estimates. Attach its projection to a minimizing disk at
$t'$ and apply Gronwall to the resulting area bound.
\end{proof}

The length identity gives a uniform constant $C_5>1$ with
\begin{equation}\label{C5eqn}
\int_{t_0}^{t_1}\int_{\widetilde\Gamma_c^\lambda(t)}k^2\,ds\,dt\le C_5.
\end{equation}

\begin{lemma}[Shi-type curvature estimate for curve-shrinking flow]
\label{lem:curve-shrinking-shi-estimate}
\source{morgan-tian-correction}
\uses{prop:curve-deformation-existence}
\label{CSSHI}


Fix $\Theta_0,L_0<\infty$. There are $\delta>0$ and constants
$\widetilde C_i$, depending only on $t_1-t_0$, an ambient curvature bound,
$\Theta_0$, and $L_0$, with the following property. If an immersed
curve-shrinking flow has initial total curvature at most $\Theta_0$ and length
at most $L_0$, and if at time $t'$ its length is at least $0<r<1$ while every
subarc of length $r$ has total curvature at most $\delta$, then, whenever
$t'+\delta r^2<t_1$,
\[
k^2\le\widetilde C_0(t-t')^{-1},\qquad
|\nabla_SH|^2\le\widetilde C_1(t-t')^{-2},\qquad
|\nabla_S^iH|^2\le\widetilde C_i(t-t')^{-(i+1)}
\]
for $t\in[t',t'+\delta r^2)$.
\end{lemma}
\begin{proof}
\sketch Lemma~\ref{lem:curve-shrinking-shi-estimate-first-inequality} proves
the curvature bound. Standard parabolic interior estimates then give the
higher derivatives; see \cite{AG,Altschuler}.
\end{proof}

\begin{remark}[The Shi-type constants depend on the initial data]
\label{rem:shi-constants-depend-on-initial-data}
\source{morgan-tian-correction}
\uses{lem:curve-shrinking-shi-estimate, prop:curve-deformation-existence}


The dependence on both $\Theta_0$ and $L_0$ is required by
Remark~\ref{rem:total-curvature-needs-initial-length}. The lifted family from
Lemma~\ref{lem:good-approximation-family} has uniform bounds for both, so the
constants are uniform in $c$ and $\lambda$.
\end{remark}

For $B>1$, Equation~\eqref{C5eqn}, Cauchy--Schwarz, and the preceding lemma
give a finite union $J_B(c,\lambda)\subset[t_0,t_1-B^{-1}]$ of a uniformly
bounded number of intervals, each of length $\delta^5B^{-2}/2$, such that its
complement has measure at most $3C_5B^{-1}$. At every time in this union,
either the ramp has length below $\delta^2B^{-1}$ or all spatial derivatives
of its curvature are uniformly bounded. After $\lambda_n\downarrow0$ and
passage to a subsequence, shrink the limiting intervals slightly to obtain
$J_B(c)$, contained in every $J_B(c,\lambda_n)$ for large $n$, with complement
of measure at most $4C_5B^{-1}$.

\begin{claim}[Piecewise domination by the comparison ODE]
\label{claim:piecewise-domination-estimate}
\uses{claim:area-forward-difference-linear-bound}
\label{C1C4delta}

Fix $C_4$ and $\widetilde A>0$. For every $\zeta>0$ there is $\delta'>0$ such
that the following holds. Let $J\subset[t_0,t_1]$ be a finite union of
intervals, let $0\le f(t_0)\le\widetilde A$, and assume
\[
f(b)\le w_{f(a),a}(b)
\]
on every component $[a,b]$ of $J$, while
$f(t'')\le f(t')+C_4(t''-t')$ for all $t'<t''$. If
$|[t_0,t_1]\setminus J|\le\delta'$, then
\[
f(t_1)\le w_{f(t_0),t_0}(t_1)+\zeta/4.
\]
\end{claim}
\begin{proof}
\uses{claim:w-a-t-explicit-formula}
\sketch Write the complementary gaps with lengths $\delta_j$ and let
$R_{\min}\ge-2C_6$. The explicit formula propagates an error through each
component of $J$ by at most $e^{C_6|J_i|}$, while each gap adds at most
$C_7\delta_j$, with $C_7=C_4+2\pi+C_6V$ and $V$ a uniform comparison-solution
bound. Thus the final error is at most
$C_7e^{C_6(t_1-t_0)}\sum_j\delta_j$.
\end{proof}

\begin{claim}[Approximate piecewise domination by the comparison ODE]
\label{claim:approximate-piecewise-domination-estimate}
\uses{claim:piecewise-domination-estimate}
\label{deltaprime}

With the preceding data and $J$ fixed, there is $\delta''>0$ such that the
same conclusion with error $\zeta/2$ holds if the component condition is
weakened to
\[
f(b)\le w_{f(a),a}(b)+\delta''.
\]
\end{claim}
\begin{proof}
\uses{claim:piecewise-domination-estimate}
\sketch Iterate the preceding error estimate. The finitely many additional errors are
products of $\delta''$ with bounded exponential factors, so choose
$\delta''$ small enough that their sum is below $\zeta/4$.
\end{proof}

\begin{claim}[Limit of the curve-shrinking flows]
\label{claim:curve-shrinking-flow-limit}
\label{limitcs}
After passing to a subsequence $\lambda_n\downarrow0$, either every large $n$
has $t_n\in J_B(c)$ with
$L(\widetilde\Gamma_c^{\lambda_n}(t_n))<\delta^2B^{-1}$, or on each component
$J_i=[t_i^-,t_i^+]$ of $J_B(c)$ the reparameterized projections
$p_1\widetilde\Gamma_c^{\lambda_n}$ converge smoothly to an immersed
curve-shrinking flow on $J_i$.
\end{claim}
\begin{proof}
\uses{claim:curve-shrinking-short-time-existence}
\sketch Outside the short-loop alternative, the curvature and all derivatives are
uniformly bounded and the lengths stay positive. Constant-speed
reparameterization and compactness give a smooth limit. Since the vertical
tangent component has integral $\lambda_n$ and bounded derivative, it tends
uniformly to zero, so the projected limit is immersed. Loss of immersion would
force curvature blow-up by the continuation criterion, contradicting the
uniform bounds.
\end{proof}

\begin{remark}[Length bound in the short-loop alternative]
\label{rem:short-loop-length-bound}
\uses{lem:immersed-curve-shrinking-area-bound, claim:curve-shrinking-flow-limit, claim:area-forward-difference-linear-bound, claim:piecewise-domination-estimate, claim:approximate-piecewise-domination-estimate, cor:curve-shrinking-length-curvature-exponential-bound, lem:curve-deformation-single-point-estimate, prop:curve-deformation-existence}
\label{xicomp}

With the later choice of $B$, the short-loop alternative supplies
$t_n\in J_B(c)$ with
\[
L(\widetilde\Gamma_c^{\lambda_n}(t_n))
<e^{-C_2(t_1-t_0)}\zeta/3.
\]
\end{remark}

\begin{lemma}[Sharper single-point estimate for the curve deformation]
\label{lem:curve-deformation-single-point-estimate}
\label{xistronger}
For every $\zeta>0$ there is $B>(t_1-t_0)^{-1}$, depending only on $\Gamma$
and the ambient curvature bound, such that, with $t_2=t_1-B^{-1}$ and
$v_c=w_{A(p_1\widetilde\Gamma(c))}$, all sufficiently small $\lambda>0$
satisfy one of:
\[
A(p_1\widetilde\Gamma_c^\lambda(t_1))<v_c(t_1)+\zeta/2,
\]
or
\[
L(\widetilde\Gamma_c^\lambda(t))<\zeta/2
\quad\text{for every }t\in[t_2,t_1].
\]
\end{lemma}
\begin{proof}
\uses{claim:area-forward-difference-linear-bound,claim:w-a-t-explicit-formula,cor:curve-shrinking-length-curvature-exponential-bound,claim:piecewise-domination-estimate,claim:curve-shrinking-short-time-existence,lem:immersed-curve-shrinking-area-bound,claim:curve-shrinking-flow-limit,rem:short-loop-length-bound,claim:approximate-piecewise-domination-estimate}
\sketch Choose $B$ so that $3C_5/B\le\delta'$, the short-loop threshold propagates to
$[t_2,t_1]$, and $e^{C_2/B}<4/3$. Apply
Claim~\ref{claim:curve-shrinking-flow-limit}. In the smooth-limit alternative,
Lemma~\ref{lem:immersed-curve-shrinking-area-bound} gives the comparison
inequality on $J_B(c)$; the linear area bound controls the complement, and
Claim~\ref{claim:approximate-piecewise-domination-estimate} gives the first
conclusion. In the other alternative, the exponential length estimate gives
the second.
\end{proof}

A finite set ${\mathcal S}\subset S^2$ is a $\nu$-net for
$\widetilde\Gamma$ if each $c$ can be joined to some $\hat c\in{\mathcal S}$
by an arc whose swept annulus has area below $\nu$; define a $\nu$-net for the
lifted family analogously. Compactness gives a common net for all sufficiently
small $\lambda$.

\begin{lemma}[Area comparison across a mu-net arc]
\label{lem:mu-net-area-comparison}
\label{muarea}
There is $\mu>0$ such that if $c,\hat c\in S^2$ are joined by an arc whose
lifted annulus at $t_0$ has area below $\mu$, then
\[
A(p_1\widetilde\Gamma_{\hat c}^\lambda(t_1))
\le v_{\hat c}(t_1)+\zeta/2
\]
implies
\[
A(p_1\widetilde\Gamma_c^\lambda(t_1))\le v_c(t_1)+\zeta.
\]
\end{lemma}
\begin{proof}
\uses{lem:minimal-annulus-area-growth}
\sketch Choose $\mu$ so annulus growth keeps the projected connecting annulus below
$\zeta/4$ at $t_1$. The explicit comparison formula also makes
$|v_c(t_1)-v_{\hat c}(t_1)|<\zeta/4$ whenever the initial disk areas differ by
at most $\mu$.
\end{proof}

\begin{lemma}[Length comparison for ramps joined by a small annulus]
\label{lem:ramp-length-comparison-small-annulus}
\uses{prop:curve-deformation-existence}
\label{BARMU}

There is $\delta>0$ such that for every $r>0$ there is $\bar\mu>0$, depending
on $r$ and the ambient curvature bound, with the following property. If ramps
$\gamma,\hat\gamma$ are joined by an annulus of area below $\bar\mu$,
$L(\gamma)\ge r$, and
\[
\int_I k\,ds<\delta
\]
for every subinterval $I\subset\gamma$ of length $r$, then
$L(\hat\gamma)\ge\frac34L(\gamma)$.
\end{lemma}
\begin{proof}
\sketch Replace an area-minimizing annulus's branch singularities using
Claim~\ref{claim:branch-point-metric-deformation}. Its smooth metric has small
area, controlled upper curvature, and the required geodesic-curvature bound on
the $\gamma$ boundary. Proposition~\ref{prop:annulus-boundary-length-estimate}
then gives the result; approximation removes boundary intersections.
\end{proof}

\begin{claim}[Length comparison across a mu-net arc]
\label{claim:mu-net-length-comparison}
\label{mulength}
For some $\mu>0$, suppose $c$ and $\hat c$ are joined by a net arc whose
lifted annulus has area at most $\mu$, and set $t_2=t_1-B^{-1}$. If
$L(\widetilde\Gamma_{\hat c}^\lambda(t))<\zeta/2$ on $[t_2,t_1]$, then
$L(p_1\widetilde\Gamma_c^\lambda(t_1))<\zeta$.
\end{claim}
\begin{proof}
  \uses{cor:curve-shrinking-length-curvature-exponential-bound, lem:ramp-length-comparison-small-annulus, lem:minimal-annulus-area-growth, lem:mu-net-area-comparison, lem:curve-deformation-single-point-estimate}

\sketch Assume the final projected length is at least $\zeta$. The exponential length
bound keeps the $c$-ramp length at least $3\zeta/4$ on $[t_2,t_1]$. Integrating
Equation~\eqref{Lformula} finds a time at which $\int k^2ds\le C_8B$; on short
arcs, Cauchy--Schwarz makes $\int kds$ smaller than the constant in
Lemma~\ref{lem:ramp-length-comparison-small-annulus}. Choose $\mu$ so annulus
growth keeps the connecting area below $\bar\mu$. That lemma then forces the
$\hat c$-ramp length to be at least $3/4$ of the $c$-ramp length, contradicting
its bound by $\zeta/2$.
\end{proof}

\begin{proposition}[Existence of the curve-shrinking deformation]
\label{prop:curve-deformation-existence}
\uses{prop:w3-forward-difference-quotient}
\label{XI}

Let $(M,g(t))$, $t_0\le t\le t_1$, be a compact Ricci flow. Given a map
$\Gamma\colon S^2\to\Lambda M$ of null-homotopic loops and $\zeta>0$, there
is a continuous family
$\widetilde\Gamma(t)\colon S^2\to\Lambda M$ of null-homotopic loops such that
$[\widetilde\Gamma(t_0)]=[\Gamma]$ in $\pi_3(M,*)$,
\[
|A(\widetilde\Gamma(t_0)(c))-A(\Gamma(c))|<\zeta,
\]
and, for every $c\in S^2$, either
$L(\widetilde\Gamma(t_1)(c))<\zeta$ or
\[
A(\widetilde\Gamma(t_1)(c))
\le w_{A(\widetilde\Gamma(t_0)(c))}(t_1)+\zeta.
\]
\end{proposition}
\begin{proof}
\uses{lem:mu-net-area-comparison,lem:curve-deformation-single-point-estimate}
\sketch Use Lemma~\ref{lem:good-approximation-family}, lift to ramps, evolve them, and
project to $M$. Choose a finite $\mu/2$-net and then one sufficiently small
$\lambda$ so Lemma~\ref{lem:curve-deformation-single-point-estimate} holds at
every net point and the lifted net has radius $\mu$. For an arbitrary $c$, the
area alternative transfers from its net point by
Lemma~\ref{lem:mu-net-area-comparison}; the short-loop alternative transfers by
Claim~\ref{claim:mu-net-length-comparison}. Projection preserves the initial
homotopy class.
\end{proof}

\section{Annuli of small area}
\label{18.5}

\begin{claim}[Smoothing the metric near interior branch points]
\label{claim:branch-point-metric-deformation}
Let $h$ be the pullback metric on an area-minimizing annulus of area at most
$\mu$, smooth away from finitely many interior branch points, with Gaussian
curvature at most $C''$ off those points. There is a smooth metric
$\widetilde h$, equal to $h$ away from small branch neighborhoods, of area at
most $2\mu$ and Gaussian curvature at most $2C''$.
\end{claim}
\begin{proof}
\sketch Near a branch point write $h=f(z,\bar z)|dz|^2$ with $f\ge0$ and
$f^{-1}(0)=\{0\}$. Replace $f$ by $f+\epsilon\rho$ for a cutoff $\rho$ that is
one near the origin. Claim~\ref{claim:perturbed-metric-curvature-bound} controls
curvature on the inner disk; smooth convergence controls the annulus where
$\rho$ changes. Choose $\epsilon$ small at each of the finitely many points to
keep the total area below $2\mu$. Boundary branch points are removable by a
smooth reparameterization \cite{White}.
\end{proof}

\begin{claim}[Curvature bound for the perturbed conformal metric]
\label{claim:perturbed-metric-curvature-bound}
If $f|dz|^2$ has Gaussian curvature at most $C''$ away from $f^{-1}(0)$, then
$(f+\epsilon)|dz|^2$ has Gaussian curvature at most $2C''$ for every
$\epsilon>0$.
\end{claim}
\begin{proof}
\sketch The curvature is
\[
\frac{-\Delta f}{2(f+\epsilon)^2}
+\frac{|\nabla f|^2}{2(f+\epsilon)^3}.
\]
From $K(f|dz|^2)\le C''$ obtain $-\Delta f\le C''f^2$ and estimate the common
numerator by $C''f^3+\epsilon C''f^2\le2C''(f+\epsilon)^3$.
\end{proof}

After scaling, it suffices to consider an annulus $A$ with $K\le1$,
$L(c_0)\ge2$, and total absolute geodesic curvature below $\delta<1/100$ on
every unit subarc of $c_0$. Choose a regular $1<\alpha<1.1$, set
$Y=\{k_{\mathrm{geod}}>\alpha\}$ and $X=\{k_{\mathrm{geod}}\le\alpha\}$. For
$x\in X$, let $D_x$ be the inward normal geodesic, stopped at length
$\delta'$ or at its first return to $\partial A$. Write $S_Z$ for the resulting
normal strip over $Z\subset X$.

\begin{remark}[Length bound on the high-curvature subset Y]
\label{rem:high-curvature-set-length-bound}
\label{Ylengthrem}
Every unit arc $J\subset c_0$ satisfies $L(J\cap Y)<\delta$.
\end{remark}

\begin{claim}[Exponential map is a local diffeomorphism near c0]
\label{claim:exponential-map-local-diffeomorphism}
\label{expcomp}
For some $0<\delta'<1/10$, the normal exponential map on $S_X$ is a local
diffeomorphism and its pullback metric is at least $(1-\delta)^2$ times the
product metric.
\end{claim}
\begin{proof}
\sketch Apply the Jacobi-field comparison determined by $K\le1$ and
$k_{\mathrm{geod}}\le\alpha$.
\end{proof}

A \emph{focusing triangle} consists of an arc $\xi\subset c_0$ with endpoints
$x,y\in X$ and initial segments of $D_x,D_y$ meeting only at their common
endpoint $v$, such that their union with $\xi$ bounds a disk $B\subset A$.
Write $a(B)=\Area(B)$ and
$t_\xi=\int_\xi k_{\mathrm{geod}}ds$.

\begin{claim}[Angle and curvature bounds for a focusing triangle]
\label{claim:focusing-triangle-angle-curvature-bound}
\label{theta/Kbd}
If $\theta_B$ is the angle at $v$, then
\[
\theta_B\le t_\xi+a(B),
\]
and, for every measurable $B'\subset B$,
\[
\theta_B-t_\xi-a(B)\le\int_{B'}K\,da<a(B).
\]
\end{claim}
\begin{proof}
\sketch Gauss--Bonnet gives
$\theta_B=t_\xi+\int_BK\,da$. Use $K^+\le1$ and split $K=K^++K^-$.
\end{proof}

\begin{claim}[Existence of a shortest geodesic to the focusing-triangle vertex]
\label{claim:focusing-triangle-vertex-geodesic}
\label{avgeo}
If $a\in\operatorname{int}\xi$ and $L(\xi)\le1$, there is a shortest
geodesic in $B$ from $a$ to $v$, meeting $\partial B$ only at its endpoints,
of length at most $1/2+\delta'$.
\end{claim}
\begin{proof}
\sketch A boundary path through the nearer endpoint has the stated length.
Compactness produces a minimizer. Any additional contact with a side admits a
corner shortcut; tangential return to $\xi$ would cut off a disk whose turning
is at most $3\pi/2+\delta$ and whose curvature integral is below $\mu$, contrary
to Gauss--Bonnet.
\end{proof}

\begin{claim}[Uniqueness of the minimal geodesic to the vertex]
\label{claim:focusing-triangle-vertex-geodesic-unique}
For every $a\in\partial A\cap B$, the minimizing geodesic in $B$ from $a$ to
$v$ is unique.
\end{claim}
\begin{proof}
\sketch Two minimizers would bound a geodesic spray. Curvature comparison makes
the exponential map locally nonsingular up to length $1/2+\delta'$, while
Gauss--Bonnet keeps all intersection angles below $\mu$. Openness and closedness
then force nearby spray geodesics to cross one minimizer, contradicting local
injectivity at its endpoint.
\end{proof}

\begin{remark}[Uniqueness of the embedded geodesic from the vertex]
\label{rem:focusing-triangle-vertex-geodesic-embedded}
The same argument gives a unique embedded geodesic in $B$ from $v$ to each
$a\in\partial A\cap B$ of length at most $1/2+\delta'$.
\end{remark}

The unique geodesics identify $B$ with a radial wedge at $v$. Write the base as
$s=h(\psi)$ and its radial contraction as $\lambda(t):s=th(\psi)$, with length
$l(t)$.

\begin{claim}[Derivative bound for the length of the radial family]
\label{claim:arc-length-derivative-bound}
\[
\frac{dl}{dt}(t)\le\max_\psi h(\psi)
\int_{\lambda(t)}k_{\mathrm{geod}}\,ds.
\]
\end{claim}
\begin{proof}
\sketch The base meets the radial sides orthogonally, so every radial
contraction does also. Its first length variation is
$\int k_{\mathrm{geod}}h|\cos\theta|ds$, which gives the estimate.
\end{proof}

\begin{lemma}[Base-length bound for a focusing triangle]
\label{lem:focusing-triangle-base-length-bound}
\label{laest}
If a focusing triangle has base $\xi$ of length at most one and bounds $B$,
then
\[
L(\xi)\le\int_\xi k_{\mathrm{geod}}\,ds+\Area(B).
\]
\end{lemma}
\begin{proof}
\uses{claim:focusing-triangle-angle-curvature-bound,claim:focusing-triangle-vertex-geodesic}
\sketch Gauss--Bonnet on each radial subwedge gives
$\int_{\lambda(t)}k_{\mathrm{geod}}ds\le t_\xi+a(B)$. Combine this with
Claim~\ref{claim:arc-length-derivative-bound},
$\max h\le1/2+\delta'$, and $l(0)=0$.
\end{proof}

\begin{corollary}[Bound on total base length of a disjoint collection of focusing triangles]
\label{cor:focusing-triangle-base-bound-collection}
\label{muepsbd}
A focusing base of length at most one has length at most $\delta+\mu$. If
focusing triangles have bases in one unit arc and disjoint interiors, the sum
of their base lengths is at most $\delta+\mu$.
\end{corollary}
\begin{proof}
\uses{lem:focusing-triangle-base-length-bound}
\sketch Sum the preceding inequality; the base turning totals at most $\delta$, and
the disjoint disks have total area at most $\mu$.
\end{proof}

\begin{lemma}[No sub-geodesic of Dx is a homotopically trivial embedded loop]
\label{lem:geodesic-no-trivial-embedded-loop}
\label{noloops}
No nonconstant sub-geodesic of any $D_x$ is an embedded null-homotopic loop in
$A$.
\end{lemma}
\begin{proof}
\sketch Such a loop has turning $\pi$ minus its closing angle. Since $K\le1$
and the enclosed area is below $\mu<\pi$, Gauss--Bonnet is impossible.
\end{proof}

For a disk $B$ cut off by some $D_x$, a unit arc $J\subset c_0\cap B$, and
$X_J=X\cap J$, stop each $D_y$, $y\in X_J$, at its first encounter with
$\partial B$ and denote the resulting strip by $S_{X_J}(B)$.

\begin{claim}[Local embedding of the geodesic family near a component of XJ]
\label{claim:local-embedding-of-geodesic-family}
\label{locemb}
For every component $C$ of $X_J$, the exponential map embeds $S_C(B)$.
\end{claim}
\begin{proof}
\sketch It is a local diffeomorphism and each radial geodesic is embedded by
Lemma~\ref{lem:geodesic-no-trivial-embedded-loop}. If two nearby radials met,
Gauss--Bonnet would bound their crossing angle by $\mu+\delta$ and force all
intermediate radials to meet one fixed radial, contradicting its embeddedness.
\end{proof}

\begin{claim}[Focusing triangles contain components of J minus XJ and admit limits]
\label{claim:focusing-triangle-limit-existence}
\label{limittri}
Every focusing triangle based in $J$ contains a component of $J\setminus X_J$.
Every infinite sequence of such triangles has a subsequence converging to a
focusing triangle.
\end{claim}
\begin{proof}
\uses{claim:local-embedding-of-geodesic-family}
\sketch The first conclusion follows from local embedding on each component of $X_J$.
Compactness gives convergent sides; their initial points remain distinct because
each base contains one of the finitely many components of $J\setminus X_J$.
\end{proof}

\begin{claim}[Disjointness of successive maximal focusing triangles]
\label{claim:maximal-focusing-triangles-disjoint}
Successive focusing triangles obtained by choosing the earliest possible left
side and the maximal base have disjoint interiors.
\end{claim}
\begin{proof}
\sketch Any crossing of successive sides either enlarges the first maximal base
or makes the right side of the first triangle the earlier eligible left side
for the second; both contradict the construction.
\end{proof}

\begin{definition}[Maximal set of focusing triangles]
\label{def:maximal-focusing-triangle-set}
The finite family obtained by repeatedly choosing the earliest possible left
side and then the maximal focusing base, removing that base before the next
choice, is the \emph{maximal set of focusing triangles} for $J$ relative to
$B$.
\end{definition}

Let $X'_J$ be $X_J$ minus all bases in this maximal family.

\begin{claim}[Length bound on the complement of the focusing-triangle bases]
\label{claim:remaining-set-length-bound}
$L(X'_J)\ge1-2\delta-\mu$.
\end{claim}
\begin{proof}
\uses{cor:focusing-triangle-base-bound-collection,rem:high-curvature-set-length-bound}
\sketch The high-curvature set removes less than $\delta$, and the disjoint focusing
bases remove at most $\delta+\mu$.
\end{proof}

\begin{claim}[Embedding of the exponential map on the remaining set]
\label{claim:exponential-map-embedding-on-remainder}
The exponential map embeds $S_{X'_J}(B)$.
\end{claim}
\begin{proof}
\sketch An intersection of radials based in $X'_J$ would form a focusing
triangle whose earliest left endpoint lies before, between, or after the
selected maximal bases. Each case contradicts the defining maximal selection.
\end{proof}

\begin{lemma}[No Dx spans between two points of c0]
\label{lem:no-geodesic-spans-c0-to-c0}
\label{Dxacross}
No $D_x$ is an embedded arc with both endpoints on $c_0$ and interior disjoint
from $\partial A$.
\end{lemma}
\begin{proof}
\uses{lem:geodesic-no-trivial-embedded-loop,claim:local-embedding-of-geodesic-family,claim:focusing-triangle-limit-existence,cor:focusing-triangle-base-bound-collection,rem:high-curvature-set-length-bound}
\sketch If $D_x$ cuts off a disk whose $c_0$-boundary has length at most one,
Gauss--Bonnet contradicts its turning and area bounds. Otherwise apply the
maximal-focusing construction to a unit subarc. Claims
\ref{claim:remaining-set-length-bound} and
\ref{claim:exponential-map-embedding-on-remainder}, together with the strip
area bound, produce a set of length at least $0.87$ whose radials end on the
disk boundary. One radial then cuts off a strictly shorter disk. Iteration
eventually reaches the forbidden length-at-most-one case.
\end{proof}

\begin{lemma}[No sub-geodesic of Dx is an embedded loop in A]
\label{lem:geodesic-no-embedded-loop-general}
No sub-geodesic of any $D_x$ is an embedded loop in $A$.
\end{lemma}
\begin{proof}
\uses{lem:no-geodesic-spans-c0-to-c0}
\sketch The null-homotopic case is
Lemma~\ref{lem:geodesic-no-trivial-embedded-loop}. For an essential loop,
compactify the complementary component containing $c_0\setminus\{x\}$ as a
disk and repeat the maximal-focusing argument. At least $0.86$ length of final
radial endpoints cannot lie on the loop segment of length below $0.2$, so some
radial joins two points of $c_0$, contradicting
Lemma~\ref{lem:no-geodesic-spans-c0-to-c0}.
\end{proof}

\begin{lemma}[Geodesics Dx, Dy from far-apart points do not meet]
\label{lem:far-apart-geodesics-disjoint}
Let initial sub-geodesics of $D_x$ and $D_y$ have a common final endpoint and
otherwise be disjoint. Their union separates $A$ with one disk component $B$;
then $B\cap c_0$ contains no arc of length one.
\end{lemma}
\begin{proof}
\uses{lem:no-geodesic-spans-c0-to-c0}
\sketch Repeat the maximal-focusing argument in $B$. The non-$c_0$ part of its boundary
has length below $0.2$, while final radial endpoints cover length at least
$0.86$, forcing a radial with both endpoints on $c_0$.
\end{proof}

\begin{corollary}[Short arcs with meeting geodesics bound a focusing disk]
\label{cor:focusing-triangle-from-short-arc}
\label{noholes}
If an arc $\xi\subset c_0$ of length at most one joins $x$ to $y$ and initial
segments of $D_x,D_y$ meet only at a common endpoint, their union with $\xi$
bounds a disk in $A$.
\end{corollary}
\begin{proof}
\uses{lem:no-geodesic-spans-c0-to-c0}
\sketch Otherwise the loop is essential; the complementary $c_0$ arc has length at
least one, contradicting Lemma~\ref{lem:far-apart-geodesics-disjoint}.
\end{proof}

\begin{corollary}[Summary of intersection properties of the geodesics Dx]
\label{cor:geodesic-family-intersection-summary}
\uses{prop:annulus-boundary-length-estimate, cor:focusing-triangle-from-short-arc, lem:no-geodesic-spans-c0-to-c0}
\label{focussummary}

Every $D_x$ is embedded and either has length $\delta'$ or ends on $c_1$. If
$D_x$ and $D_{x'}$ meet, the points are joined on $c_0$ by an arc $\xi$ of
length at most $\delta+\mu$; suitable initial segments and $\xi$ bound a disk
$B$, and
\[
L(\xi)\le\int_\xi k_{\mathrm{geod}}\,ds+\Area(B).
\]
\end{corollary}
\begin{proof}
\sketch Combine Lemmas~\ref{lem:geodesic-no-embedded-loop-general} and
\ref{lem:far-apart-geodesics-disjoint} with
Corollary~\ref{cor:focusing-triangle-from-short-arc} and
Lemma~\ref{lem:focusing-triangle-base-length-bound}.
\end{proof}

\begin{claim}[Total base length of focusing triangles meeting a unit interval]
\label{claim:focusing-triangle-total-base-length-bound}
For every unit subarc $J\subset c_0$, the total length of maximal focusing
bases meeting $J$ is at most $2\delta+\mu<0.03$.
\end{claim}
\begin{proof}
\uses{cor:focusing-triangle-base-bound-collection}
\sketch Each base has length at most $\delta+\mu$, so those meeting $J$ lie in an arc
of length below two. Their total turning is at most $2\delta$ and the disjoint
triangle areas total at most $\mu$.
\end{proof}

\begin{proposition}[Boundary length estimate for annuli of small area]
\label{prop:annulus-boundary-length-estimate}
\label{annlengthest}
Fix $0<\delta<1/100$. For every $r>0$ and $C''<\infty$ there is $\mu>0$ such
that the following holds. Let $A$ be an annulus with boundary $c_0\sqcup c_1$
and Gaussian curvature at most $C''$. If $L(c_0)>r$, every length-$r$ subarc
$I\subset c_0$ satisfies $\int_I|k_{\mathrm{geod}}|ds<\delta$, and
$\Area(A)<\mu$, then
\[
L(c_1)\ge\frac34L(c_0).
\]
\end{proposition}
\begin{proof}
\sketch Scale to $r=1$, $K\le1$, and $L(c_0)\ge2$. Cut the annulus along one
radial joining $c_0$ to $c_1$, and remove the maximal focusing bases. By
Claim~\ref{claim:focusing-triangle-total-base-length-bound}, these remove less
than $0.06L(c_0)$; the high-curvature set and the strip-area estimate remove
less than another $0.10L(c_0)$. On the remaining set $X''$, of length at least
$0.84L(c_0)$, the exponential map is embedded and every radial reaches $c_1$.
Claim~\ref{claim:exponential-map-local-diffeomorphism} makes the terminal arc
length at least $(1-\delta)L(X'')>0.83L(c_0)>3L(c_0)/4$.
\end{proof}

\section{Local curvature estimate for curve-shrinking flow}
\label{18.6}

Fix a curve-shrinking flow $c(x,t)$ in a Ricci flow $(W,h(t))$. At a time
$t'$, assume the curve has length at least $r$ and every length-$r$ subarc has
total curvature at most $\delta$. Let $t_2$ be maximal such that
$k^2\le2/(t-t')$ on $[t',t_2]$.

\begin{claim}[Uniform bound on length and total curvature on short sub-arcs]
\label{claim:curve-shrinking-length-curvature-uniform-bound}
\label{C0claim}
There is $D_0$, depending only on the initial total-curvature bound and ambient
curvature bound, such that every initial subarc $\gamma_{t'}$ of length at most
$r$ satisfies
\[
\int_{\gamma_t}k\,ds<D_0,
\qquad L(\gamma_t)\le D_0r
\]
on $[t',t'+\delta r^2]$.
\end{claim}
\begin{proof}
\uses{cor:curve-shrinking-length-curvature-exponential-bound}
\sketch Apply the exponential length and total-curvature bound to the whole curve and
restrict to the evolved subarc.
\end{proof}

For an initial length-$r$ arc centered at $x_0$, let
$s(x,t)=\int_{x_0}^x|X(y,t)|dy$ and choose a fixed cutoff $\psi$ supported in
$[-1/2,1/2]$, equal to one on $[-3/8,3/8]$. Set
$\varphi(x,t)=\psi(s(x,t)/r)$.

\begin{claim}[Bound on the time derivative of the cutoff function]
\label{claim:cutoff-function-time-derivative-bound}
\label{1887}
There is $D_1$, depending only on ambient curvature, such that
\[
|\partial_t\varphi(x,t)|
\le\frac{D_1}{r\sqrt{t-t'}}+D_1
\]
on $[t',t_2]$.
\end{claim}
\begin{proof}
\uses{lem:curve-shrinking-first-variation,claim:curve-shrinking-length-curvature-uniform-bound}
\sketch Equation~\eqref{dLdteqn} bounds
$|\partial_ts|$ by an ambient-curvature multiple of $L(\gamma_t)$ plus
$\int k^2ds$. Use $k^2\le\sqrt{2/(t-t')}k$, the preceding claim, and
$\partial_t\varphi=r^{-1}\psi'\partial_ts$.
\end{proof}

\begin{claim}[Bound on the cutoff-weighted curvature integral]
\label{claim:cutoff-weighted-curvature-integral-bound}
\label{claim1885}
There are $D_2$, depending on the initial total-curvature and ambient bounds,
and $D_3$, depending only on ambient curvature, such that
\[
\left|\frac d{dt}\int_{\gamma_t}\varphi k\,ds\right|
\le D_2\left(1+\frac1{r\sqrt{t-t'}}\right)+\frac{D_2}{r^2}
+D_3\int_{\gamma_t}\varphi k\,ds.
\]
\end{claim}
\begin{proof}
\uses{claim:cutoff-function-time-derivative-bound,claim:curve-shrinking-curvature-inequality,lem:curve-shrinking-length-curvature-growth,claim:curve-shrinking-length-curvature-uniform-bound}
\sketch Differentiate the weighted integral. The cutoff-time term is controlled by the
preceding claim. The regularized curvature inequality reduces the other term
to $\int\varphi k''ds+C\int\varphi kds$; two integrations by parts and
$|\varphi''|\le C/r^2$ give the displayed estimate.
\end{proof}

\begin{corollary}[Square-root bound on the cutoff-weighted curvature integral]
\label{cor:cutoff-weighted-curvature-integral-sqrt-bound}
For $t\in[t',t_2]$,
\[
\int_{\gamma_t}\varphi k\,ds\le D_4\sqrt\delta,
\]
where $D_4$ depends only on the initial total-curvature and ambient bounds.
\end{corollary}
\begin{proof}
\sketch Integrate Claim~\ref{claim:cutoff-weighted-curvature-integral-bound}
over an interval of length at most $\delta r^2$, using
$\int_{\gamma_{t'}}kds<\delta$, $r<1$, and Gronwall.
\end{proof}

\begin{corollary}[Uniform bound on the curvature integral over short sub-arcs]
\label{cor:curvature-integral-uniform-bound}
\label{D_4cor}
Every evolved initial subarc of length at most $r$, and every time-$t$ subarc
of length at most $r/2$, satisfies
\[
\int k\,ds\le2D_4\sqrt\delta
\]
for $t\in[t',t_2]$.
\end{corollary}
\begin{proof}
\uses{cor:curve-shrinking-length-curvature-exponential-bound}
\sketch Cover an initial arc by the central regions of two cutoff arcs. The lower
length comparison transports any time-$t$ arc of length $r/2$ into an initial
arc of length at most $r$.
\end{proof}

Use exponential coordinates at the midpoint of an initial arc, with Euclidean
metric $h'$ and $z$-axis tangent to the curve. For sufficiently small $r_0$,
all $h(t)$ are $C^1$-close to $h'$ on the coordinate ball. The curvature bound
moves points by at most $4\sqrt{t-t'}\le4\sqrt\delta r$. On a central interval
$I'$ of $h'$-length between $0.8r$ and $r$, write the evolving curve as a graph
$(z,f(z,t))$ while $|f_z|_{h'}<1/10$. Let $Z=(1,f_z)$ and let $Y$ be the
graph-flow velocity.

\begin{equation}\label{1814}
\nabla_SS=\frac{\nabla_ZZ}{|Z|^2}
-\frac{\langle\nabla_ZZ,Z\rangle}{|Z|^4}Z.
\end{equation}

\begin{claim}[Comparison of the graph-flow and curve-shrinking velocities]
\label{claim:graph-flow-velocity-comparison}
\label{curvflowcompare}
With
$\psi(Z)=\langle\Gamma(Z,Z),\partial_z\rangle_{h'}/|Z|^2$,
\[
Y=\frac{\nabla_ZZ-\langle\Gamma(Z,Z),\partial_z\rangle_{h'}Z}{|Z|^2}
=\nabla_SS+left(
\frac{\langle\nabla_ZZ,Z\rangle}{|Z|^4}-\psi(Z)
\right)Z.
\]
\end{claim}
\begin{proof}
\sketch In Euclidean coordinates
$\nabla_ZZ=(0,f_{zz})+\Gamma(Z,Z)$. The displayed vector is $h'$-orthogonal to
the $z$-axis and differs from Equation~\eqref{1814} by a multiple of $Z$;
these two properties characterize graph-flow velocity.
\end{proof}

\begin{claim}[Angle and norm bounds for the graph flow]
\label{claim:graph-flow-angle-and-norm-bounds}
\label{1891}
For sufficiently small $\delta$, the angle between $Y$ and $Z$ exceeds
$\pi/4$, and
\begin{enumerate}
\item $k\le|Y|<\sqrt2k$;
\item $|\langle\nabla_ZZ,Z\rangle|<(k+2\delta)|Z|^3$;
\item $|\nabla_ZZ|<2(|Y|+\delta)$;
\item $|\langle Y,Z\rangle|\le|Y||f_z|(1+3\delta)$.
\end{enumerate}
\end{claim}
\begin{proof}
\sketch Since $|f_z|<1/10$, $(0,f_{zz})$ makes a uniformly positive angle with
$Z$. Metric closeness transfers this to $h(t)$ and compares $Y$ with its
normal component $\nabla_SS$. The remaining estimates follow from
Claim~\ref{claim:graph-flow-velocity-comparison}, $|\psi|<1.5\delta$, and
$h'$-orthogonality of $Y$ to the $z$-axis.
\end{proof}

\begin{claim}[Bounds for the connection correction term]
\label{claim:graph-flow-connection-term-bounds}
\label{psiclaim}
For an ambient-curvature constant $C'$,
\[
|Z(\psi(Z))|<C'(1+\delta|Y|),
\qquad
|Y(\psi(Z))|<C'(|Y|+\delta|\nabla_ZY|).
\]
\end{claim}
\begin{proof}
\sketch Differentiate the definition of $\psi$. Derivatives of $\Gamma$ are
bounded by ambient curvature, and derivatives landing on $Z$ are controlled by
$\delta|\nabla_ZZ|$ or $\delta|\nabla_YZ|$. Use the preceding claim and
$[Y,Z]=0$.
\end{proof}

The graph-flow length element satisfies
\begin{equation}\label{Z_t}
\partial_t|Z|^2
=Z\left(\frac{Z(|Z|^2)}{|Z|^2}\right)-2|Y|^2|Z|^2+V,
\end{equation}
where
\begin{equation}\label{Vest}
|V|<C'(1+\delta|Y|).
\end{equation}

\begin{claim}[Mixed derivative bound for the graph flow]
\label{claim:graph-flow-mixed-derivative-bound}
\[
|\langle\nabla_ZY,Z\rangle|
\le\bigl(|f_z|(|\nabla_ZY|+2\delta|Y|)+\delta|Z|^2|Y|\bigr)(1+\delta).
\]
\end{claim}
\begin{proof}
\sketch Write $Y=(0,\phi)$ in graph coordinates. Then
$\nabla_ZY=(0,\phi_z)+\Gamma(Z,Y)$, so its $h'$ inner product with
$Z=(1,f_z)$ is bounded by $|f_z||\phi_z|+\delta|Z|^2|Y|$. Compare $h'$ and
$h(t)$ and use $|\phi_z|\le|\nabla_ZY|+\delta|Z||Y|$.
\end{proof}

\begin{lemma}[First curvature bound in the curve-shrinking Shi estimate]
\label{lem:curve-shrinking-shi-estimate-first-inequality}
\label{csshi1}
Let $(W,h(t))$, $t_0\le t\le t_1$, be a Ricci flow and fix
$\Theta<\infty$. There are $\delta>0$ and $0<r_0\le1$, depending only on
$\Theta$ and the ambient curvature bound, such that if an immersed
curve-shrinking flow has initial total curvature at most $\Theta$, and at some
$t'$ and $0<r\le r_0$ has length at least $r$ and total curvature at most
$\delta$ on every length-$r$ subarc, then, provided
$t'\le t_1-\delta r^2$,
\[
k^2\le\frac2{t-t'}
\qquad\text{on }[t',t'+\delta r^2].
\]
\end{lemma}
\begin{proof}
\sketch On the maximal interval $[t',t_2]$ where the asserted estimate holds,
the cutoff estimates make the total curvature of every short arc
$O(\sqrt\delta)$. Hence each central arc remains a graph. Integrating
Equation~\eqref{Z_t} over $I'$ gives an $L^2$ bound for $f_z$; together with
the total-curvature bound, this improves $|f_z|<1/10$ and shows that the graph
description persists to $t_2$.

For the graph flow, set
\[
Q=\frac{|Y|^2}{2-|Z|^2}.
\]
Equations~\eqref{Z_t}--\eqref{Vest} and
Claims~\ref{claim:graph-flow-angle-and-norm-bounds}--\ref{claim:graph-flow-mixed-derivative-bound}
give, after absorbing the $\delta$ errors,
\[
Q_t\le\frac{Q_{zz}}{|Z|^2}-(Q-C_1')^2+(C_1')^2.
\]
If $l=L_{h'}(I')$, compare $Q-C_1'$ with
\[
\frac1{t-t'}+rac{4(1-\delta)^{-1}l^2}{z^2(l-z)^2}+C_1'.
\]
The maximum principle and $l\ge0.8r$ imply at the midpoint
$Q<3/(2(t-t'))$. Since $|Z|^2\ge1-\delta$ and $k\le|Y|$, this improves the
bootstrap to $k^2<2/(t-t')$. Thus $t_2=t'+\delta r^2$. The graph maximum
principle follows the local argument of \cite{AG}; the corresponding graphical
hypersurface estimate is in \cite{EckerHuisken}.
\end{proof}
